\documentclass[11pt]{article}
\usepackage[margin=0.9in]{geometry}
\usepackage[T1]{fontenc}
\usepackage{lmodern,microtype,amsmath,amssymb,booktabs,graphicx,array}
\usepackage[colorlinks=true,linkcolor=blue,citecolor=blue,urlcolor=blue]{hyperref}
\hypersetup{pdftitle={Graph Geometry and Input-Entry Effects in SHA-256: A Controlled Measurement Study},pdfauthor={Bee Rosa Davis},pdfkeywords={SHA-256, spectral graphs, heat kernel, avalanche effect, higher-order differentials}}
\usepackage[font=small,labelfont=bf]{caption}
\usepackage[section]{placeins}
\setlength{\parskip}{4pt}
\setlength{\parindent}{0pt}
\setlength{\emergencystretch}{2em}
\setcounter{topnumber}{4}
\setcounter{totalnumber}{5}
\renewcommand{\topfraction}{0.9}
\renewcommand{\textfraction}{0.07}
\renewcommand{\floatpagefraction}{0.8}
\newcommand{\HoldoutDifference}{+0.009}
\newcommand{\HoldoutCI}{[-0.2869, 0.3140]}
\newcommand{\SearchDifference}{+0.126}
\newcommand{\SearchCI}{[-0.4354, 0.6657]}
\newcommand{\FinalHamming}{128.024}
\newcommand{\RFAccuracy}{50.315}
\newcommand{\SVMAccuracy}{49.800}
\newcommand{\GraphRows}{%
Primary & 1.0 & Silhouette & 0.8957 & 0.9066 & -0.0110 [-0.0439, 0.0196] \\
Primary & 1.0 & Gap & 0.2050 & 0.2046 & +0.0003 [-0.0349, 0.0422] \\
Primary & 8.0 & Silhouette & 0.1608 & 0.1635 & -0.0028 [-0.0047, -0.0010] \\
Primary & 8.0 & Gap & 0.6159 & 0.6139 & +0.0020 [-0.0009, 0.0051] \\
Primary & -- & Detour & 2.5629 & 2.5660 & -0.0031 [-0.0053, -0.0008] \\
Confirmation & 1.0 & Silhouette & 0.8803 & 0.8985 & -0.0182 [-0.0471, 0.0059] \\
Confirmation & 1.0 & Gap & 0.2461 & 0.2246 & +0.0215 [-0.0145, 0.0599] \\
Confirmation & 8.0 & Silhouette & 0.1654 & 0.1647 & +0.0007 [-0.0030, 0.0049] \\
Confirmation & 8.0 & Gap & 0.6145 & 0.6153 & -0.0008 [-0.0032, 0.0017] \\
Confirmation & -- & Detour & 2.5640 & 2.5674 & -0.0034 [-0.0065, -0.0006] \\%
}
\newcommand{\SensitivityRows}{%
state -- PCG64 & -0.00045 & [-0.0032, 0.0025] \\
state -- MT19937 & +0.00042 & [-0.0026, 0.0033] \\
digest -- PCG64 & +0.00092 & [-0.0027, 0.0045] \\
digest -- MT19937 & +0.00179 & [-0.0010, 0.0044] \\
PCG64 -- MT19937 & +0.00087 & [-0.0019, 0.0036] \\%
}
\newcommand{\CubeRows}{%
8 & 48.83 & 49.52 & -0.69 & [-5.21, 3.68] & 1.000 \\
10 & 46.48 & 50.58 & -4.09 & [-8.74, 0.36] & 0.975 \\
12 & 53.12 & 49.07 & +4.05 & [-0.41, 8.36] & 1.000 \\
14 & 47.07 & 49.20 & -2.13 & [-6.48, 2.30] & 1.000 \\
16 & 51.56 & 50.06 & +1.50 & [-3.10, 5.97] & 1.000 \\
18 & 49.80 & 49.59 & +0.21 & [-4.19, 4.46] & 1.000 \\
20 & 47.46 & 50.52 & -3.06 & [-7.38, 1.41] & 1.000 \\
24 & 52.54 & 49.94 & +2.60 & [-1.75, 7.01] & 1.000 \\%
}
\newcommand{\MLRows}{%
RF & 50.315 & 49.825--51.050 & 0.5017 \\
PCA + SVM & 49.800 & 48.225--51.425 & 0.5003 \\%
}

\title{Graph Geometry and Input-Entry Effects in SHA-256:\\A Controlled Measurement Study}
\author{Bee Rosa Davis\\Davis Geometric\\\small \href{mailto:bee_davis@alumni.brown.edu}{bee\_davis@alumni.brown.edu}\\\small ORCID: \href{https://orcid.org/0009-0009-8034-4308}{0009-0009-8034-4308}}
\date{Working paper}
\begin{document}
\maketitle
\begin{abstract}
Geometric descriptions of cryptographic state spaces depend on how states are sampled, connected and perturbed. We study these dependencies in SHA-256 using matched random-bit graphs, coordinate-resolved avalanche measurements, Boolean cube sums and held-out digest classifiers. Two groups of twenty independent SHA--random dataset pairs show that high spectral-clustering silhouette and graph detour occur in both populations and change substantially with graph construction. Small detour contrasts appear in two groups; a third, generator-sensitivity sample does not resolve the contrast. Across 112,640 one-bit message interventions, positions selected for lower discovery-sample Hamming distance do not retain that ordering on independent bases: the held-out selected-minus-other difference is \HoldoutDifference{} bits. Equal-budget two-bit searches and six-bit cubes do not establish an advantage for the fixed candidate set tested. Input-word entry and exact register transport account for prominent timing differences in a separate cube control. Random-forest and PCA--SVM classifiers attain mean held-out digest accuracies of \RFAccuracy\% and \SVMAccuracy\%. The contribution is a reproducible measurement protocol that distinguishes graph representation, intervention timing and tested discrimination performance; it is not a security proof or a collision/preimage attack.
\end{abstract}
\noindent\textbf{Keywords:} SHA-256; spectral graph analysis; heat kernel; avalanche effect; higher-order differentials; reproducibility.

\section{Question and scope}
SHA-256 is a deterministic computation with visible algebraic organization, even when its outputs are useful as cryptographic hashes. Finding organization in intermediate states is therefore a different question from constructing a distinguisher or reducing the cost of a cryptanalytic task. A measurement must identify both the observed object and the comparison that makes its result informative.

We investigate three questions. First, do graph statistics of terminal working states differ from statistics of independently sampled bit strings subjected to the same graph construction? Second, do input coordinates selected by a geometric-distance criterion retain an advantage on new messages and under equal intervention budgets? Third, how much of an apparent leading-versus-trailing register separation follows from input timing and register transport?

The analysis concerns a fixed set of representations and tests. It does not assert that SHA states have no detectable structure, that no other selection rule can help an attack, or that failure of a tested classifier establishes security. The central empirical distinction is between a statistic that is large in absolute terms and a statistic that discriminates under a matched comparison.

\section{Primitives, coordinates and validation}
We implement the SHA-256 compression steps specified in FIPS 180-4~\cite{nist}. Let $X_r(m)=(a_r,b_r,c_r,d_r,e_r,f_r,g_r,h_r)$ be the eight working words after $r$ completed rounds of the first compression block. $X_0$ is the initial value; $X_{64}$ precedes feed-forward. The final digest $H(m)$ includes feed-forward and all padding blocks. For a single block and fixed initial value, feed-forward is a fixed word-wise modular-addition bijection, but it need not preserve bit distances or a fitted classifier's features.

Message-based experiments use independent 55-byte inputs, so SHA padding produces one block. An input coordinate $i$ is MSB-first within a byte: byte $\lfloor i/8\rfloor$, mask $2^{7-(i\bmod8)}$. The corresponding message word is $W_{\lfloor i/32\rfloor}$. These are input coordinates, not coordinates of working register $a$ or any other register. All 440 content bits are eligible for the avalanche experiment; the last content word has 24 eligible bits.

The scalar implementation was compared with \texttt{hashlib} at padding boundaries, including messages that require multiple blocks. A vectorized implementation of raw-block working states was independently checked against the scalar captures at rounds 0, 1, 2, 8, 15, 16, 17, 18, 32 and 64. Block index and within-block round are recorded separately. Regression gates cover positive-degree graph normalization, complete heat traces, truncation bounds, reference calibration, bit placement and GF(2) rank. The release includes the tests and their results.

\begin{table}[ht]
\centering\small
\caption{Experimental units. Shared interventions on one base are not counted as independent base messages.}
\begin{tabular}{@{}p{.25\linewidth}p{.66\linewidth}@{}}\toprule
Experiment & Design \\\midrule
Spectral graphs & Two groups of 20 independent dataset pairs; 512 states per dataset; two kernel parameters; full spectra. \\
Detour sensitivity & 20 fresh datasets per representation: working state, digest, PCG64 and MT19937 bit strings. \\
Avalanche & 128 discovery and 128 holdout bases; all 440 single-bit interventions on each; 27 sampled rounds. \\
Two-bit search & 256 new bases; candidate set and 20 random sets; 15 two-bit patterns per set and base. \\
Six-bit cubes & 512 new bases; candidate set and 20 random sets; 64 corners; eight round counts. \\
Input-entry control & 256 uniform raw blocks; three-bit cubes placed in five message words; eight relative round counts. \\
Digest classifiers & Five independent fits per model; 4,000 training and 4,000 test observations per fit, balanced by class. \\\bottomrule
\end{tabular}
\end{table}

\section{Graph construction and matched reference populations}
Binary working states are embedded in $\mathbb R^{256}$ without fitted dimensional reduction. Euclidean distance is the square root of Hamming distance. For each point, we find 30 nonself nearest neighbors and form directed weights
\begin{equation}
 W^{\rightarrow}_{ij}=\mathbf 1\{j\in N_{30}(i)\}
 \exp\!\left(-\frac{\|x_i-x_j\|^2}{4h}\right),\qquad
 W=\frac{W^{\rightarrow}+(W^{\rightarrow})^\top}{2}.
\end{equation}
We study $h=1$ and $h=8$ with identical settings for SHA states and independently sampled random bits. With $d_i=\sum_jW_{ij}>0$, the symmetric normalized Laplacian is
\begin{equation}L=I-D^{-1/2}WD^{-1/2}.\end{equation}
Its zero eigenvector is proportional to $\sqrt d$; a common positive rescaling of $W$ leaves $L$ unchanged. These invariants are checked numerically. Nearest-neighbor construction, bandwidth and normalization are components of the estimator, not properties inferred from the data~\cite{luxburg}.

We compute the complete 512-mode eigensystem. KMeans, with five clusters and ten initializations, is fitted to the five eigenvectors immediately after the zero mode. We report the resulting silhouette on those coordinates. The reference dataset receives the same fitted pipeline. Random labels on SHA coordinates would answer a different question.

A second statistic, the \emph{detour ratio}, uses an undirected union graph of ten nonself neighbors, with Euclidean edge lengths. It is the mean shortest-path distance divided by direct Euclidean distance over finite, distinct pairs. Sparse graphs can have ratios well above one on random clouds. Thus the random-data comparator is measured, not assigned the value one.

\subsection{Heat observables and interpretation}
For graph diffusion time $t\geq0$, distinct from compression-round count $r$,
\begin{equation}
 K_t=\exp(-tL),\quad Z(t)=\operatorname{tr}K_t=\sum_{j=0}^{n-1}e^{-t\lambda_j},\quad
 K_t(i,i)=\sum_j e^{-t\lambda_j}\phi_j(i)^2.
\end{equation}
The release saves complete spectra and traces at $t=0.01,0.1,1,10$. In particular, $Z(0)=n$. A truncated computation retaining $k$ smallest modes instead has $Z_k(0)=k$ and satisfies $0\leq Z(t)-Z_k(t)\leq(n-k)e^{-t\lambda_{k-1}}$. Full spectra avoid that truncation issue in this study.

The finite-graph diagonal is nonincreasing because its derivative is $-\sum_j\lambda_j e^{-t\lambda_j}\phi_j(i)^2\leq0$. A two-time diagonal ratio therefore does not establish the sign of an underlying manifold's scalar curvature. Likewise, a Poincare representation with prescribed negative curvature imposes a metric. When binary sign vectors are projected to a common radius, the resulting distance is a monotone function of Hamming distance. Neither construction alone demonstrates an intrinsic hyperbolic manifold.

\subsection{Replication and results}
Each primary replication uses a new seed and independently generated SHA inputs and random bits. We use percentile bootstrap intervals with 5,000 resamples of the paired dataset-level differences. These are descriptive, unadjusted intervals. They do not treat graph edges or eigenmodes as independent observations.

\begin{figure}[ht]\centering\includegraphics[width=\linewidth]{figures/graph_controls.pdf}
\caption{Each dot is one dataset's fitted silhouette; short bars show means. The same pipeline produces high scores in SHA and random data at $h=1$. Broadening the kernel changes the score substantially in both populations. The confirmation group uses fresh seeds.}\end{figure}

\begin{table}[ht]\centering\small
\footnotesize
\caption{Graph statistics, with SHA-minus-random differences and descriptive 95\% bootstrap intervals. Every row has 20 paired dataset replications.}
\begin{tabular}{@{}lllr r l@{}}\toprule
Group & $h$ & Statistic & SHA & Random & Difference [95\% interval] \\\midrule
\GraphRows
\bottomrule\end{tabular}\end{table}

The kernel parameter changes mean silhouette from approximately 0.90 to 0.16 in both populations. The absolute height of a silhouette score therefore does not identify a SHA-specific effect. The primary $h=8$ silhouette difference is slightly negative; that contrast does not repeat in the confirmation group. The spectral-gap intervals include zero in both groups at both kernel parameters.

The detour contrast is more nuanced. Its small negative interval repeats in the first two groups. We therefore retain that observation rather than replacing it with a claim of equality. A further twenty-dataset sensitivity sample compares working states and actual digests with both PCG64 and MT19937 references (Table~\ref{tab:sensitivity}). Those intervals include zero, and the estimates vary with representation and reference. These observations leave a small construction-specific detour contrast unresolved; they do not establish either a robust general fingerprint or its absence.

\begin{table}[ht]\centering\small
\caption{Fresh detour sensitivity sample, 20 paired replications per comparison. State means terminal pre-feed-forward working state. These follow-up comparisons were specified after observing the first two graph groups.}\label{tab:sensitivity}
\begin{tabular}{@{}lrr@{}}\toprule
Comparison & Mean difference & 95\% interval \\\midrule
\SensitivityRows\bottomrule
\end{tabular}\end{table}

The confirmation and sensitivity samples are follow-ups, not a preregistered external replication. Their protocols and seeds were saved before their execution. Reporting each group separately preserves the distinction between initial exploration and fresh evaluation without assuming that a later nonsignificant interval disproves an earlier effect.

\section{Coordinate-resolved avalanche and selection}
For each base message $m$ and input bit $i$, we measure
\begin{equation}A_r(m,i)=\operatorname{HW}\bigl(X_r(m)\oplus X_r(m\oplus e_i)\bigr),\end{equation}
where $\operatorname{HW}$ counts changed bits across all eight working words. This is a physical input intervention with a declared output statistic. It is not a per-output-bit diffusion rate.

Every one of 128 discovery bases is perturbed at all 440 content positions. The 22 positions with smallest discovery mean $A_{24}$ are frozen before evaluating a separate set of 128 bases. Selecting 22 positions defines a candidate set; the count is not itself a test for exceptional bits. Resampling is over bases, keeping all within-base interventions together.

\begin{figure}[ht]\centering\includegraphics[width=\linewidth]{figures/avalanche_entry.pdf}
\caption{Measured holdout avalanche. Left: mean Hamming distance grouped by the actual flipped message word. Right: selected word groups aligned by their entry round. Later words begin influencing the state later; a full-trajectory slope would combine that delay with growth and saturation.}\end{figure}

The selected-minus-other holdout difference at round 24 is \HoldoutDifference{} bits, with interval \HoldoutCI{}. The data do not establish a persistent lower-distance advantage for this selection rule. Across the holdout interventions, mean round-64 Hamming distance is \FinalHamming{} bits out of 256. This is an observed average, not a claim that all directions, messages or higher-order statistics are identical.

We separately test the fixed candidate positions $S=\{1,5,7,9,10,14\}$ on 256 fresh bases. For each six-position set, the search evaluates all 15 two-position flips and retains the smallest raw round-24 Hamming distance. The comparator comprises twenty uniformly sampled six-position subsets of W0, with exactly the same search budget. Distances are left on the same scale on both sides.

\begin{figure}[ht]\centering\includegraphics[width=\linewidth]{figures/holdout_search.pdf}
\caption{Independent evaluation of selected coordinates. Left: discovery versus holdout means for all 440 positions, with the frozen 22 highlighted. Right: per-base candidate-minus-control difference for equally budgeted two-bit searches; negative values favor the candidate.}\end{figure}

The candidate-minus-mean-control difference is \SearchDifference{} bits, with interval \SearchCI{}. This comparison does not resolve an advantage for the tested candidate. The interval is conditional on the twenty sampled control sets; it does not quantify uncertainty over every possible six-position set. The search question remains distinct from digest discrimination or attack complexity.

\section{Cube measurements and the timing of interventions}
For a declared set of $k$ input coordinates $S$, the cube sum of an output $f$ is its $k$-fold Boolean derivative,
\begin{equation}\Delta_S f(m)=\bigoplus_{u\in\{0,1\}^{k}}f(m\oplus E_Su).\end{equation}
A Boolean polynomial of total degree below $k$ has zero $k$-fold derivative. A zero sum on a particular cube does not, conversely, prove a global degree bound. Cube techniques are established tools for testing polynomial structure~\cite{dinur}; here we report explicitly bounded measurements of working-state projections.

\subsection{Candidate cubes with matched entry time}
We compare the candidate set $S$ with twenty random six-position subsets of W0. All variables therefore enter at the same compression step. The same 512 independently generated 55-byte bases are used for every set. Each cube evaluates 64 corners. The observable is the XOR sum of bit 16, numbered from the least-significant end, of register $a$. A zero bit-sum has reference probability one half under the random-function comparison.

\begin{table}[ht]\centering\small
\caption{Six-bit cube results. Differences and intervals are percentage points. The last column adjusts the eight candidate-versus-one-half binomial tests by Holm's method; it is not a test of the candidate-versus-control difference.}
\begin{tabular}{@{}rrrrlr@{}}\toprule
Round & Candidate (\%) & Control (\%) & Difference & 95\% interval & Holm $p$ \\\midrule
\CubeRows\bottomrule
\end{tabular}\end{table}

None of the eight candidate zero-rate tests rejects after the stated adjustment. Every paired candidate-minus-control interval also includes zero. This supports a bounded negative result for this set, projection, background distribution and sample size. It does not show equivalence to all controls or bound the best cube over all coordinates and observables.

\subsection{Input-entry and register transport}
An input variable in original message word $W_j$ first enters at compression step $j$, or completed-round count $j+1$. We use $\delta=r-(j+1)$ to describe time since that entry. A separate control draws 256 uniform raw 64-byte blocks and moves the same three MSB-first offsets $\{0,1,2\}$ among W0, W4, W8, W12 and W15. These are raw compression-block interventions, including positions that would be padding or length bits in a message-based experiment; no padded-message interpretation is applied.

SHA-256 also copies working registers between rounds:
\begin{equation}
 b_r=a_{r-1},\quad c_r=a_{r-2},\quad d_r=a_{r-3},\qquad
 f_r=e_{r-1},\quad g_r=e_{r-2},\quad h_r=e_{r-3}.
\end{equation}
For the same base and intervention, cube differentiation commutes with these identities. Thus $\Delta_S d_r=\Delta_S a_{r-3}$ and $\Delta_S h_r=\Delta_S e_{r-3}$ whenever the displayed histories are defined. A trailing register can retain a zero sum after a leading register changes through exact transport, without an additional geometric hypothesis.

\begin{figure}[ht]\centering\includegraphics[width=\linewidth]{figures/cube_controls.pdf}
\caption{Two distinct cube observables. Left and center: fraction of raw bases with a whole-word zero sum for the three-bit entry control, aligned to variable entry; W0 and W15 illustrate the transport pattern. Right: single-output-bit zero-rate differences for the six-bit candidate experiment, with base-resampled 95\% intervals. A whole-word zero and a zero bit have different reference probabilities.}\end{figure}

The entry control displays the leading-to-trailing sequence predicted by the register identities. Its exact zero fractions depend on the selected offsets and backgrounds, so we do not identify a universal round threshold. Total round count alone is an incomplete description of the experiment: both the entry time of the variables and the observation register are necessary to interpret a surviving cube sum.

\section{Held-out classification of actual digests}
The classification experiment uses $H(m)$ rather than $X_{64}(m)$. Each of five independent replications contains 2,000 SHA digests and 2,000 independent random-bit strings for training, and another balanced 4,000-observation holdout. The fixed models are a 100-tree random forest with maximum depth ten and an RBF SVM with $C=1$, default scaled gamma and fifty PCA components. PCA is fitted on training observations only. No model or feature selection uses the holdout.

\begin{table}[ht]\centering\small
\caption{Five independent classifier replications. Ranges describe the five fits and are not confidence bounds on all distinguishers.}
\begin{tabular}{@{}lrrr@{}}\toprule
Model & Mean accuracy (\%) & Accuracy range (\%) & Mean AUC \\\midrule
\MLRows\bottomrule\end{tabular}\end{table}

\begin{figure}[ht]\centering\includegraphics[width=\linewidth]{figures/digest_classifiers.pdf}
\caption{Each dot is a model trained and evaluated on a fresh balanced dataset. Horizontal colored bars indicate means; the dotted line is one-half. These results quantify the tested procedures, rather than cryptographic indistinguishability.}\end{figure}

The observed means are close to chance, with individual fits varying on either side. This experiment has finite sample size and limited model classes. It supplies no basis for inferring full-round security from classifier failure, nor does it invalidate a different statistic evaluated on a different object.

\section{Implications and reproducibility}
The measurements separate three sources of apparent organization. Graph construction can make both SHA and random points highly clusterable in a chosen spectral coordinate system. Input placement creates a delay before an intervention can affect any register. Register transport then carries a changed leading word into trailing positions over subsequent rounds. A useful state-space analysis must control all three before attaching a cryptanalytic interpretation to an observed pattern.

The results leave open further work on small graph-statistic contrasts and alternative intervention rules. Stronger claims would require independent seeds and reference generators, bandwidth and sample-size sensitivity, a frozen selection rule, and an explicitly defined attack or prediction task. Our data already show why absolute clustering scores, selected quantiles, and unqualified round counts are insufficient summaries.

The accompanying \texttt{sha256\_reproducibility.zip} supplies the exact main-run source snapshot, follow-up runners, protocols, seeds, dependency versions, full graph spectra, intervention arrays, cube indicators, classifier holdouts and predictions, figure builder and manuscript source. Every numerical table and figure is generated from those records. No external market or proprietary input data are required. A repository DOI can be added when this archive is deposited; no public deposit identifier is asserted in this version.

All bootstrap intervals use independently sampled bases or whole dataset pairs as the resampling unit. Most are descriptive and unadjusted; only the stated eight cube binomial tests use a familywise adjustment. The shared candidate/control bases and the sampled control sets are explicitly retained. The graph follow-ups were chosen after initial results and are reported as such. These design boundaries delimit the findings without converting the paper into a claim about every possible SHA-256 measurement.

\paragraph{Funding.} This research received no external funding.
\paragraph{Author statement.} The author develops the geometric research framework motivating these measurements. This study evaluates the specified computational procedures and does not establish the broader framework's mathematical claims.

\begin{thebibliography}{9}
\bibitem{nist} National Institute of Standards and Technology. \emph{Secure Hash Standard (SHS)}. FIPS PUB 180-4, 2015. \url{https://doi.org/10.6028/NIST.FIPS.180-4}.
\bibitem{luxburg} U. von Luxburg. A tutorial on spectral clustering. \emph{Statistics and Computing}, 17:395--416, 2007. \url{https://arxiv.org/abs/0711.0189}.
\bibitem{dinur} I. Dinur and A. Shamir. Cube attacks on tweakable black box polynomials. \emph{EUROCRYPT 2009}, pp. 278--299. Preprint: \url{https://eprint.iacr.org/2008/385}.
\end{thebibliography}
\end{document}
